\section{A complete local strong-II Gram with independent testing rates}
\label{sec:r7-local-strong-II}

\begin{theorem}[Strong-II local allocation]
Let the positive C2 source satisfy $z,y\in\Inf_{>0}$ and let
$B=H+L$ be a positive unit.  Keep the exact source coefficients of
(GE1), with $n\geq0$ and the remaining displayed weights positive.
Then the original C2 parent belongs to $\sum(K_c\otimes_{\mathbb R}
K_t)^2$ whenever both independent affine testing coordinates $u,v$
are infinitesimal, including either literal testing axis.
There is no comparison hypothesis between $u^2,z$, between $v^2,y$,
or between the two source rates.  All openings of $r$, $L/y$,
$1+zr-L/y$, and $1-H$ are retained, as is literal $n=0$.
\end{theorem}

This theorem is a full-parent local result.  The remote testing charts
and the original projective endpoints are separate obligations; they
must be joined before the strong II/HH checklist row is admitted.
All identities below are polynomial in testing coordinates.  Inverses
are taken only in the pure source field, and the final step uses the
bounded Gram strict-real-residue lemma already used in the completed
local UI certificate.  In particular no ordering or inverse of the
mixed tensor is used.

\paragraph{Source units and notation.}
Write $U=u^2,V=v^2,k=\sqrt y$ and
\[
 a_c=1+zr,\quad t=L/y,\quad j=a_c-t,\quad s=a_c/B,
 \quad c=CL,\quad d=Dyj.
\]
The source entry gives bounded $t,j,r$, positive units $B,c,d,m/z^2,A/y$,
bounded $nr$ and $\alpha(1-H)$, and strong $\alpha$.
These facts include the arbitrary openings $t\to0$, $j\to0$,
$r\to0$, and $1-H\to0$.  They follow directly from the positive
source formulas and were proved in (DI7).  Let $h=\operatorname{st}B$;
then $0<h\leq1$ and, with $r_0=\operatorname{st}r$,
\[
\begin{gathered}
 c_0=d_0^{\rm pair}=\frac{h}{2(1+h)},\quad
 \operatorname{st}(m/z^2)=\frac{1-r_0}{2h^2(1+h)},\quad
 \operatorname{st}(A/y)=\frac1{2h(1+h)},\\
 \operatorname{st}(nr)=\frac1{2(1+h)}-\frac{r_0}{2},\qquad
 \operatorname{st}\{\alpha(1-H)\}=\frac{h^2}{2(1+h)}.
                                                        \tag{LII1}
\end{gathered}
\]
These residues follow directly from (GE1): $L\to0$, $H\to h$,
$a_c\to1$, $G\to h$, $\delta\to h(1+h)$, and the source
numerator $M\to h^2$.  Substitution in the six rational coefficient
formulas gives (LII1), including the identity
$\operatorname{st}(nr)=1/[2(1+h)]-r_0/2$ without division by $n$.
Define the source quantities (with $q,g$ bounded)
\[
 \ell=y(C+D)=c/t+d/j,\qquad
 q=\frac{D}{C+D}=\frac{dt}{cj+dt},\qquad
 g=\frac{yCD}{C+D}=\frac{cd}{cj+dt}.
\]
Here $0<q<1$, $g$ is a positive unit with residue $h/[2(1+h)]$,
and $\ell$ has a positive real lower bound, although it may be
unbounded.  Indeed $cj+dt$ is a unit because $t+j=a_c$ is a unit.
Choose positive real $\ell_*\leq\ell$ and
$\alpha_*\leq\alpha$.

\paragraph{The state roots and the exact C,D completion.}
Use the source-fixed parity graph
\[
 X=-\frac{zr v}{B}W+\xi,\qquad
 Y=\frac{Hu}{B}W+\eta,\qquad
 Z=\frac{uv}{B}W+\zeta.                         \tag{LII2}
\]
It agrees with the actual kernel at $U=z,V=y$ and is a change of all
four homogeneous state coordinates with determinant one.  Set
\[
 Q_n=u\xi-\zeta,\quad K=ukW,\quad
 R_h=\frac{vQ_n-su(V-y)W}{k},\qquad
 N=\frac\eta k-qR_h.                              \tag{LII3}
\]
Thus $\eta=k(N+qR_h)$.  Every root in (LII3) vanishes at the actual
contact.  If $Q=uW-Y$ and $J=Q+v(uX-Z)$, then
\[
 Q/k=(t/B)K-N-qR_h,\qquad
 J/k=(1-q)R_h-N-(j/B)K.
\]
Consequently the exact weighted completion is
\[
\begin{split}
 4CQ^2+4DJ^2={}&4\ell N^2+4gR_h^2
 +\frac{8(d-c)}BKN-8gsKR_h
 +\frac{4(ct+dj)}{B^2}K^2,\\
 4CVQ^2={}&4CV\eta^2
 +\frac{4ct}{B^2}VK^2-\frac{8c}BVK(N+qR_h).
                                                        \end{split}\tag{LII4}
\]
There is no estimate in these identities.  Retain the following
four nonnegative, contact-centered packets outside the remainder:
\[
 4(\ell-\ell_*)N^2,\qquad 4CV\eta^2,\qquad
 4n(1+U)(1+V)\xi^2,\qquad
 4(\alpha-\alpha_*)(1+U)(v\eta-\zeta)^2.          \tag{LII5}
\]
Call the exact remainder $G$.  Its coefficients are bounded after
writing $\eta=k(N+qR_h)$ and using (LII4).  For completeness, its
entire formula is
\[
\begin{split}
G={}&4\ell_*N^2+4gR_h^2+\frac{8(d-c)}BKN-8gsKR_h\\
 &+\frac{4(ct+dj)}{B^2}K^2
 +\frac{4ct}{B^2}VK^2-\frac{8c}BVK(N+qR_h)\\
 &+4n(1+U)(1+V)
   \left\{\frac{z^2r^2V}{B^2}W^2-\frac{2zr v}B W\xi\right\}\\
 &+4\alpha_*(1+U)(v\eta-\zeta)^2
  +4\alpha(1+U)
   \left\{\frac{(1-H)^2UV}{B^2}W^2
        -\frac{2(1-H)uv}B W(v\eta-\zeta)\right\}\\
 &+4mM^2+4A(1+U)(1+V)Y^2\\
 &+2\{J^2+(1+U)(vY-Z)^2+(1+V)(X+uZ)^2
                      -(1+U)(1+V)Z^2\},        \end{split}\tag{LII6}
\]
where $M=vW-X+uvY-uZ$ and (LII2) is used in the last line.
Every potentially unbounded original coefficient occurs here only
through $nr$, $\alpha(1-H)$, or the bounded quantities in (LII1).
No division by a testing coordinate has been introduced.

\paragraph{One exact gauge removes the dependent-root first normal.}
In the formal bounded polynomial (LII6), the coefficient of $R_hW$
is $uk\rho(U,V)$ for a bounded source polynomial $\rho$.  Replace
this entire term using the exact identity
\[
 KR_h=uvWQ_n-sU(V-y)W^2.                         \tag{LII7}
\]
Call the resulting polynomial $\widetilde G$.  It is equal to $G$
on the roots (LII3), and has no $R_hW$ term.  This gauge is useful
because $R_h$ is a dependent root: it permits the state and angular
contact equations to be read as formal coefficient vanishings.

More precisely, by parity write uniquely
\[
 \widetilde G=Q_0(N,\xi,\zeta,R_h)
 +W\{vP_\xi(U,V)\xi+uvP_\zeta(U,V)\zeta
                         +ukP_N(U,V)N\}
 +\phi(U,V)W^2,                                  \tag{LII8}
\]
where all coefficients are bounded source polynomials and $Q_0$ is
quadratic in the four displayed roots.  Formula (LII6), followed by
the one replacement (LII7), is an exact coefficient-extraction
specification for every entry in (LII8); the accompanying checker
also prints them without suppression.

The retained packets (LII5) vanish together with their first jets at
the actual contact.  The exact original source state equations thus
give
\[
 P_\xi(z,y)=P_\zeta(z,y)=P_N(z,y)=0.              \tag{LII9}
\]
Indeed the physical $\eta$ derivative first forces the $N$ coefficient
to vanish; the physical $\xi,\zeta$ derivatives then force their
formal coefficients to vanish because the $R_hW$ coefficient was
removed.  The exact original contact and the two angular derivatives
now give
\[
 \phi(z,y)=\phi_U(z,y)=\phi_V(z,y)=0.              \tag{LII10}
\]
This uses the six exact original source substitutions already checked
in R6.  It imposes no new relation between source and testing fields.

\paragraph{Bounded division and the complete angular allocation.}
There are two additional elementary support facts in (LII8):
\[
 P_\xi\in(U,z),\qquad \deg_U\phi\leq2,\qquad
 \phi(0,V)=V(d_s+e_sV),                           \tag{LII11}
\]
where
\[
 d_s=\frac4{B^2}\{m(B+zr)^2+(n+\tfrac12)z^2r^2\},\qquad
 e_s=\frac4{B^2}(n+\tfrac12)z^2r^2.
\]
Both $d_s/z^2$ and $e_s/z^2$ are bounded.  The first support assertion
is literal bounded-coefficient membership: setting $U=z=0$ in the
extracted polynomial gives zero.  It is also clear before expansion:
all $\xi W$ terms not supplied by the gauge contain $U$, $m$, $nzr$,
or $zr$; those supplied by the gauge contain $U$.

Monic division using (LII9) therefore gives bounded polynomials
\[
\begin{split}
 P_\xi&=(U-z)A_\xi+z(V-y)B_\xi,\\
 P_\zeta&=(U-z)A_\zeta+(V-y)B_\zeta,\\
 P_N&=(U-z)A_N+(V-y)B_N.                         \end{split}\tag{LII12}
\]
For the first line, evaluate $P_\xi$ at $U=z$: the support assertion
supplies the factor $z$ with bounded quotient, and its value at $V=y$
is zero.  All other divisions are by the monic polynomials $U-z$
or $V-y$.  In particular none creates an inverse of a small source
or a testing difference.

The scalar has the exact bounded decomposition
\[
 \phi=\frac{d_s+e_sV}{z^2}V(U-z)^2
       +U(U-z)(V-y)A_s(V)+U(V-y)^2C_s(V).        \tag{LII13}
\]
Here is a constructive proof, including boundedness.  Write
$\phi=a(V)U^2+b(V)U+V(d_s+e_sV)$ and subtract the first term of
(LII13).  The remaining $U^2,U$ coefficients $a_1,b_1$ vanish at
$V=y$ by the value and $U$ derivative in (LII10).  Divide them by
$V-y$ to obtain $A_s,B_s$.  The $V$ derivative gives
$zA_s(y)+B_s(y)=0$, so set $C_s=(B_s+zA_s)/(V-y)$.
All three divisions are monic divisions of bounded polynomials.
This proves (LII13) exactly, not asymptotically.

Introduce the two legal angular roots
\[
 R=v(U-z)W,\qquad T=uk(U-z)W.                    \tag{LII14}
\]
The following identities allocate every term in (LII12)--(LII13):
\[
\begin{split}
 zv(V-y)W&=\frac{uVQ_n-uvkR_h}{s}-(V-y)R,\\
 uv(V-y)W&=\frac{VQ_n-vkR_h}{s},\\
 uk(V-y)W&=\frac{vkQ_n-yR_h}{s},\\
 U(U-z)(V-y)W^2&=-\frac{R_hT-uRQ_n}{s},\\
 U(V-y)^2W^2&=\frac{VQ_n^2-2vkR_hQ_n+yR_h^2}{s^2}.
                                                        \end{split}\tag{LII15}
\]
Together with $v(U-z)W=R$, $uv(U-z)W=uR$, and
$uk(U-z)W=T$, these are an explicit complete allocation of the
full remainder as a bounded symmetric Gram in
\[
                    (N,\xi,\zeta,R_h,R,T).       \tag{LII16}
\]
The original first state normals are all present.  The scalar
restriction has not been used as a substitute for them.

\paragraph{The finite block and the missing pivot.}
The residue of $Q_0$ on the first four roots is
\[
             \operatorname{diag}(4\ell_*,2,4\alpha_*,4g_0).
                                                        \tag{LII17}
\]
For example the signed core at $W=0$ is exactly
$2(\xi+u\zeta-uv\eta)^2+2(1+U)V\xi^2+2(1+V)\eta^2$;
all its other entries vanish in the present residue.  The extra
terms in the first three lines of (LII15) have infinitesimal Gram
coefficients except for $A_\xi\xi R$ and the allowed $A_NNT$.
Likewise the last line has only infinitesimal coefficients, and
the fourth line has one potentially finite $R_hT$ entry.  Thus,
apart from entries in the last row and column $T$, the sole finite
off-diagonal entry is between $\xi$ and $R$.

The exact coefficient extraction yields
\[
 \operatorname{st}A_\xi=\frac4{h(1+h)},\qquad
 \operatorname{st}(d_s/z^2)=\frac2{h^2(1+h)}.     \tag{LII18}
\]
For a short independent check of the first value, the unscaled
original first-normal coefficient on the $V=y$ slice is
\[
 -4nzr(1+y)+4m(1-H)-(4D+2)yj+2(1+y)
 \longrightarrow 2-4d_0^{\rm pair}=\frac2{1+h}.
\]
Its full multiplier is $2/B$.  Passing from the original normal
to (LII8) changes this $U$-quotient only by a bounded multiple of
$U$ or $z$, so its residue is unchanged.  The second value follows
from (LII1):
\[
 \operatorname{st}\{m(B+zr)^2/z^2+(n+\tfrac12)r^2\}
 =\frac{1-r_0}{2(1+h)}+\frac{r_0}{2(1+h)}
 =\frac1{2(1+h)}.
\]
Hence the finite $(\xi,R)$ Gram block is exactly
\[
 \begin{pmatrix}
  2&\dfrac2{h(1+h)}\\[2mm]
  \dfrac2{h(1+h)}&\dfrac2{h^2(1+h)}
 \end{pmatrix},\qquad
 \text{Schur complement}=\frac2{h(1+h)^2}>0.      \tag{LII19}
\]
This calculation retains $r_0=0$, $n=0$, and $h=1$.

Finally the exact toric conductor is the zero identity
\[
 \boxed{\quad
 T^2-UR^2-\frac{U-z}{s}R_hT
             +\frac{u(U-z)}s RQ_n=0.\quad}       \tag{LII20}
\]
Add any positive real multiple $\Lambda$ of this zero identity to
the Gram.  Its only nonzero residue contribution is $\Lambda T^2$.
The five-by-five block excluding $T$ is strictly positive by
(LII17)--(LII19); all entries incident to $T$ are bounded.  Choose
a sufficiently large real $\Lambda$ by the ordinary real Schur
complement.  The entire six-by-six residue Gram is then strictly
positive.  The bounded Gram lemma gives a tensor SOS for $G$, and
adding the four packets (LII5) proves the theorem.

\paragraph{Quantifiers and endpoint discipline.}
No proof step divides by $u$, $v$, $U-z$, or $V-y$: parity quotients
and monic polynomial divisions are coefficient extractions whose
final identities are polynomial.  Thus the two literal affine axes
are included.  The source inverses $1/B$, $1/k$, and the source
packet square roots are legal even when their values are large;
all coefficients of the final Gram are bounded.  The source ratios
$t,j,r,1-H$ have not been bounded away from zero.  The new roots
raise testing degree, so the local certificate is not homogenized
to assert the original literal projective endpoints.  Those are
joined by the separate direct-parent remote certificate.

\paragraph{Checkable scope.}
\texttt{check\_local\_strong\_ii.py} checks (LII4), (LII7),
(LII15), (LII20), the complete bounded remainder's parity and support,
its scalar $U=0$ value, its hard residue, and all three values in
(LII18)--(LII19).  It prints the exact extracted polynomials in
\texttt{extracted\_coefficients.txt}.  Source contact and stationarity
are inherited from the six exact original source checks in
\texttt{analysis\_atlas/oct19\_r6\_c2\_double/check\_common\_ii.py};
no additional generic source substitution is claimed here.
